\chapter{Automorphic Forms on a Quaternion Algebra}

\label{ch:10}

Let $F$ denote a global field and $D$ a (division) quaternion algebra defined over $F$. Thus $D$ is a central simple (division) algebra of rank 4 over $F$. The multiplicative group $G'$ of $D$ is a reductive algebraic group defined over $F$ and the homogeneous space

\[
\bX' = Z_\infty^+ \bG'_F \backslash \bG'_A
\]

possesses a $\bG'_A$-invariant measure.

By an cusp form on $\bG'_A$ we shall understand any irreducible unitary representation of $\bG_A$ which occurs in the decomposition of the natural representation $R'$ of $\bG_A$ in $\bL^2(\bX')$. The purpose of this section then is to describe the proof and consequences of the following result:

\begin{theorem}\label{thm:10.1}

There is a correspondence which associates to each irreducible unitary representation $\pi'$ of $\bG'_A$ an irreducible unitary representation $\pi$ of $\GL(2, \A)$ so that if $\pi'$ is a cusp form on $\bG'_A$ then $\pi$ is a cusp form on $\GL(2, \A)$ (provided $\pi'$ is not one-dimensional).

\end{theorem}

This theorem does for quaternion algebras what Theorem~\ref{thm:7.11} does for quadratic extensions of $F$. Both theorems provide natural solutions to the problem (Problem A) of relating automorphic forms on some group $G'$ to automorphic forms on $\GL(2)$.

To establish the correspondence alluded to in the first part of Theorem~\ref{thm:10.1} one applies the construction of Weil with $V$ equal to $D$ instead of a quadratic extension $L$ (cf.\ Sections~\ref{ch:7}.A and B). To prove the second (and more difficult) part of the theorem one could follow the lead of Theorem~\ref{thm:7.11} and analyze the $L$-function $L(s, \pi)$ attached to $\pi$. In this case

\begin{equation}\label{eq:10.1}
L(s, \pi) = L(s, \pi')\tag{10.1}
\end{equation}

where $L(s, \pi')$ is the $L$-function naturally attached to $\pi'$ via the method of Tate alluded to in the ``Notes and References'' of Section~\ref{ch:6}. Since this $L$-function can be shown to be entire and to satisfy the requisite functional equation (again by Tate's method) Theorem~\ref{thm:10.1} follows from Jacquet-Langlands' characterization of cusp forms on $\GL(2)$.

The proof just sketched (in particular the identity \eqref{eq:10.1}) demonstrates the naturality of the correspondence $\pi' \to \pi$. Nevertheless in Subsection B below we present a different proof which seems to be the ``best possible.'' This proof involves the Selberg trace formula and is natural in the sense that it establishes an essentially functional analytic result (the occurrence of $\pi$ in a certain $\bL^2$ space) using only functional analysis (not the intermediary device of $L$-functions). This proof also allows one to significantly strengthen Theorem~\ref{thm:10.1} as follows:

\begin{theorem}\label{thm:10.2}

Let $S$ denote the finite set of places $v$ in $F$ such that $D_v = D \otimes_F F_v$ is (still) a division algebra. Then the correspondence

\[
\pi' \to \pi = \otimes \pi_v
\]

is one-to-one onto the collection of cusp forms on $\GL(2, \A(F))$ with $\pi_v$ square-integrable for each $v \in S$.

\end{theorem}

In Subsection C we shall explain some of the possible consequences of this remarkable theorem. To appreciate one particular application we have in mind it is helpful to understand the connection between automorphic forms on a quaternion algebra and the classical theory of cusp forms. The point is that thus far we have restricted attention to automorphic forms on $\GL(2)$ and hence ultimately to classical forms automorphic with respect to some arithmetic Fuchsian group of congruence type. By dealing with division algebras one is led to consider certain arithmetic groups with compact fundamental domain in $\SL(2, \R)$. Automorphic forms with respect to such $\Gamma$ do not possess familiar Fourier expansions (because $\Gamma$ has no unipotent elements) and consequently such forms have vigorously resisted explicit construction and systematic analysis.

Several preliminary results are collected in Subsection A below. In the Subsection (D) we sketch yet another proof of Theorem~\ref{thm:10.1} (due to Shimizu). This proof directly constructs the space of (theta) functions on which $\pi$ acts. Therefore it is analogous to Shalika-Tanaka's treatment of Theorem~\ref{thm:7.11}. Its approach is desirable because it says something very definite about the representability of classical cusp forms by theta-series associated to quaternary quadratic forms (a long-standing problem).

\section{Preliminaries}\label{sec:10a}

Let $D$ denote a division quaternion algebra defined over $F$ and for each place $v$ in $F$ put

\[
D_v = D \otimes_F F_v.
\]

We say then that $v$ is ramified in $D$ if $D_v$ is a division algebra. Otherwise $v$ is unramified and $D_v$ is naturally isomorphic to $M(2, F_v)$. (We also say that $D$ splits (or does not split at $v$) if $v$ is unramified (resp.\ ramified) in $D$.)

The number of $F$-places ramified in $D$ is finite and even and the set of such places will be denoted by $S$. This set completely determines $D$, i.e.\ if there is given a set consisting of an even number of places of $F$, then there exists a unique quaternion algebra $D$ over $F$ such that $S$ is precisely the set of places in $F$ ramified in $D$. (Unique up to isomorphism.)

For each place $v$, put

\[
\bG'_v = D_v^{\times}.
\]

Following the theory for $\GL(2)$ one should attach to each local group $\bG'_v$ a Hecke group algebra $\mathcal{H}(\bG'_v)$.

If $v$ is unramified in $D$ there is an isomorphism

\[
\theta_v: D_v \to M(2, F_v).
\]

To fix such an isomorphism we first fix a maximal order $\mathfrak{P}$ in $D$. (An order in $D$ is any subring which contains a basis for $D$ over $F$ and is finite as a $\mathcal{O}_F$-module.) Letting $\mathfrak{P}_v$ be the $\mathcal{O}_v$-module in $D_v$ generated by $\mathfrak{P}$ we fix $\theta_v$ such that $\theta_v(\mathfrak{P}_v) = M(2, \mathcal{O}_v)$. With this choice of $\theta_v$ we may identify $\bG'_v$ with $\GL(2, F_v)$ and define $\mathcal{H}(\bG'_v)$ as the inverse image (under $\theta_v$) of the Hecke algebra of $\GL(2, F_v)$.

We define $\bK'_v$ to be the maximal compact subgroup of $\bG'_v$ such that $\theta_v(\bK'_v) = \bK_v = \GL(2, \mathcal{O}_v)$. Thus for unramified $v$ the notion of an admissible representation of $\mathcal{H}(\bG'_v)$ is obvious.

If $v$ is ramified in $D$ then $\mathfrak{P}_v$ is the unique maximal order in $D_v$, namely $\{x \in D_v: |N(x)|_v \leq 1\}$. Thus

\[
\bK'_v = \{x \in \bG'_v: |N(x)|_v = 1\}
\]

is the maximal compact subgroup of $\bG'_v$ and $\mathcal{H}(\bG'_v)$ is defined to be the algebra of locally constant compactly supported functions on $\bG'_v$. (If $v$ is archimedean, $\bK_v$ is simply the norm one group of Hamilton's quaternion algebra, and $\mathcal{H}(\bG'_v)$ is the sum of the space of measures defined by infinitely differentiable compactly supported functions on $\bG'_v$ which are bi-$\bK'_v$-finite and the space of measures defined by matrix coefficients of finite dimensional representations of $\bK'_v$.)

An admissible representation $\pi'$ of $\mathcal{H}(\bG'_v)$ for $v$ ramified has (once again) the obvious definition. For such an irreducible representation there exists an irreducible representation of $\bG'_v$ which ``lifts'' to $\pi'$ on $\mathcal{H}(\bG'_v)$. This representation of $\bG'_v$ will be admissible in the sense of Definition~\ref{def:4.9}. Actually it will be finite dimensional since $\bG'_v$ is compact modulo its center. Conversely, every finite-dimensional continuous representation of $\bG'_v$ is admissible.

As for the global theory, let

\[
\bK' = \prod_v \bK'_v.
\]

Let $\mathcal{H}(\bG'_A)$ denote the space spanned by all functions $f' = \prod f'_v$ with $f'_v \in \mathcal{H}(\bG'_v)$ and $f'_v$ equal to the characteristic function of $\bK'_v$ for almost every $v$. Then a representation $\pi'$ of $\mathcal{H}(\bG'_A)$ in $V$ is called admissible if:

\begin{enumerate}[(i)]

\item Every $w$ in $V$ is of the form $\sum_i \pi'(f_i)w_i$ with $f_i$ a function in $\mathcal{H}(\bG'_A)$;

\item $\pi'(\xi)$ is finite-dimensional if $\xi$ is an elementary idempotent of $\mathcal{H}(\bG'_A)$ (i.e.\ of the form $\prod \xi_v$ with $\xi_v$ a linear combination of characters of irreducible representations of $\bK'_v$).

\end{enumerate}

As for $\GL(2)$, a representation of $\bG'_A$ will be called admissible if its restriction to $\bK'$ contains any irreducible of $\bK'$ at most finitely many times. It should not seem surprising that every irreducible such representation defines such a representation of $\mathcal{H}(\bG'_A)$ and is completely factorizable as

\[
\bigotimes_v \pi'_v.
\]

\subsection*{Some Arithmetic}

In this paragraph we collect some facts which make it possible to relate the homogeneous space $Z_\infty^+ \bG'_F \backslash \bG'_A$ to the classical space

\[
\Gamma \backslash \{\operatorname{Im}(z) > 0\}.
\]

Here $\Gamma$ is some lattice subgroup of $\SL(2, \R)$. In particular we assume that $F$ is a totally real number field and that $D$ is indefinite over $F$, i.e.\ at least one infinite place of $F$ is unramified in $D$. To simplify matters as much as possible we shall actually assume that $F = \Q$.

Much of the arithmetic quoted here has been developed in greater generality by Eichler. (See the references at the end of this section.)

Let $\mathfrak{P}$ denote any maximal order in $D$. Suppose $\mathfrak{A}$ is a finite $\Z$-module such that $\Q\mathfrak{A} = D$. Then $\mathfrak{A}$ is a right (resp.\ left) $\mathfrak{P}$-ideal if $\mathfrak{A}\mathfrak{P} \subset \mathfrak{A}$ (resp.\ $\mathfrak{P}\mathfrak{A} \subset \mathfrak{A}$). Two right (resp.\ left) $\mathfrak{P}$-ideals $\mathfrak{A}$ and $\mathfrak{B}$ are called equivalent if $\delta\mathfrak{A} = \mathfrak{B}$ (resp.\ $\mathfrak{A}\delta = \mathfrak{B}$) for some $\delta \in D$. The number of equivalence classes of right $\mathfrak{P}$-ideals is finite and equal to the number of equivalence classes of left ideals. This number is independent of the choice of maximal order $\mathfrak{P}$ and (by analogy with the case of fields) is called the class number of $D$. For maximal orders and indefinite $D$ the class number is one.

Now consider the group of ideles $\bG'_A$ of $D$. Its subgroup of $x$ such that $x_p$ is a unit of $\mathfrak{P}_p$ for all finite $p$ is precisely $\bG'_\infty \prod_{p < \infty} \bK'_p$. On the other hand, for every $x \in \bG'_A$, there is a (uniquely determined) right $\mathfrak{P}$-ideal $\mathfrak{A}$ such that $\mathfrak{A}_p = x_p\mathfrak{P}_p$. If we denote this right $\mathfrak{P}$-ideal $\mathfrak{A}$ by $\cap x_p\mathfrak{P}_p$ then the mapping

\[
\mathbf{x} \to \cap_p x_p\mathfrak{P}_p
\]

establishes a one-to-one correspondence between $D^{\times}\backslash \bG'_A/\bG'_\infty \prod_{p < \infty} \bK'_p$ and equivalence classes of right $\mathfrak{P}$-ideals.

\begin{lemma}\label{lem:10.3}

$\bG'_A = \bG'_\Q(\bG'_\infty)^{0}\prod_{p < \infty}\bK'_p$.

\end{lemma}

\begin{proof}

Since $\bG'_\infty = \GL(2, \R)$, $(\bG'_\infty)^{0} = \GL^{+}(2, \R)$. From the remarks above (in particular the fact that the class number of $D$ is one) it follows that $\bG'_A = \bG'_\Q \bG'_\infty \prod \bK'_p$. Thus the lemma follows from the fact that maximal orders contain units of norm $-1$.

\end{proof}

\begin{remark*}

If $F$ is any number field then

\[
\A(F)^{\times} = F^{\times}(F_\infty)^{0}\left(\prod_{v\ \text{finite}} \mathcal{O}_v^{\times}\right)
\]

if and only if the class number of $F$ is one. That is, there is a one-to-one correspondence between equivalence classes of $\mathcal{O}$-ideals (right or left since $F$ is commutative) and the homogeneous space

\[
F^{\times}\backslash \A(F)^{\times}/F_\infty^{0}\left(\prod_{v\ \text{finite}} \mathcal{O}_v^{\times}\right).
\]

(Compare the ``Notes and References'' of Section~\ref{ch:3}.) So Lemma~\ref{lem:10.3} simply generalizes this fact to non-abelian division algebras and its proof too follows the abelian case.

\end{remark*}

\begin{corollary*}

\[
Z_\infty^+ \bG'_\Q \backslash \bG'_A / \prod_{p < \infty} \bK'_p \cong \Gamma \backslash \SL(2, \R)\tag{10.2}
\]

where

\[
\Gamma = \bG'_\Q \cap (\bG'_\infty)^+ \prod_{p < \infty} \bK'_p.\tag{10.3}
\]

\end{corollary*}

Thus the norm one group of a maximal order of a quaternion division algebra provides a Fuchsian group with compact fundamental domain.

The compactness of the quotient space

\[
Z_\infty^+ \bG'_F \backslash \bG'_A
\]

is established, for example, in \cite{Weil4}.

\begin{remark}\label{rem:10.4}

The dense subspace of $\bK'$-finite $\mathfrak{z}'$-finite functions in $\bL^2(\bX')$ coincides with the space of functions $\phi(g)$ on $\bG'_A$ which satisfy the following properties:

\begin{enumerate}[(i)]

\item $\phi(\gamma g) = \phi(g)$ for all $\gamma \in \bG'_F$;

\item $\phi(zg) = \phi(g)$ for all $z \in Z_\infty^+$;

\item $\phi(g)$ is right $\bK'$-finite;

\item as a function of $\bG'_\infty$ alone, $\phi$ is right $\mathfrak{z}'$-finite ($\mathfrak{z}'$ being the center of the universal enveloping algebra of the complexified Lie algebra of $\bG'_\infty$); and

\item $\phi$ is ``slowly increasing.''

\end{enumerate}

We denote this space by $A(\bX')$. Since $\bX'$ is compact, the continuity of $\phi(g)$ (together with the transforming property (ii) above) actually implies that Condition (v) is redundant. We include it to stress that $A(\bX')$ is precisely the analogue for $D$ of the space of automorphic forms for $\GL(2)$.

\end{remark}

In case $D$ is an indefinite division algebra over $\Q$, $\bG'_\infty = \GL(2, \R)$. In this case Lemma~\ref{lem:10.3} implies that special functions in $A(\bX')$ (namely those satisfying \eqref{eq:3.6}(iii), \eqref{eq:3.6}(iv) and $\phi(gk') = \phi(g)$ for all $k \in \prod_{p < \infty} \bK'_p$) correspond to functions $f(z)$ in $\{\operatorname{Im}(z) > 0\}$ satisfying the following properties;

\[
\begin{aligned}
\text{(i)} \quad &f\!\left(\frac{az + b}{cz + d}\right) = (cz + d)^{k}f(z) \text{ for all } \begin{pmatrix} a & b \\ c & d \end{pmatrix} \in \Gamma = \bG'_\Q \cap \GL^{+}(2, \R)\\
& \prod_{p < \infty} \bK'_p; \text{ and }
\end{aligned}
\]

(ii) $f(z)$ is holomorphic.

Hence the terminology ``automorphic form'' for an element of $A(\bX')$ (or for an irreducible constituent of $R'_\psi$) is apt. Note that no cuspidal or boundedness condition is relevant since

\[
\Gamma \backslash \{\operatorname{Im}(z) > 0\}
\]

is compact. Note also that if $\phi$ is right $\SO(2)$-invariant and $\phi$ satisfies $\Delta\phi = \frac{1 - s^2}{4}\phi$ (in place of Conditions \eqref{eq:3.6}(iii) and (iv)) then $f_\phi(z)$ is real analytic and a wave-form in the sense of Maass.

\section{Statement and Proof of the Fundamental Result}\label{sec:10b}

According to Theorem~\ref{thm:7.6} one may use Weil's representation to associate to each irreducible unitary representation $\pi'$ of a local division quaternion algebra (over $F$) an irreducible unitary representation $\pi(\pi')$ of $\GL(2, F)$. This representation $\pi$ will be a special representation if $\pi'$ is one-dimensional and super cuspidal otherwise. Furthermore all such representations of $\GL(2, F)$ are so obtained.

\begin{theorem}\label{thm:10.5}

Let $D$ denote a division quaternion algebra defined over the global field $F$, $S$ the set of places of $F$ ramified in $D$, and $G'$ the multiplicative group of $D$. To each irreducible unitary admissible representation $\pi' = \bigotimes_v \pi'_v$ of $\bG'_A$ let $\pi$ denote the representation of $\bG_A = \GL(2, \A)$ whose $v$-th component is equivalent to $\pi'_v$ if $v \notin S$ and equivalent to $\pi_v(\pi'_v)$ if $v \in S$. (Here $\pi_v(\pi'_v)$ denotes the irreducible component of the Weil representation $r(D_v)$ indexed by $\pi'_v$.) THEN:

\begin{enumerate}[(i)]

\item $\pi = \otimes \pi_v$ is a cusp form for $\bG_A$ if $\pi'$ is a (greater than one dimensional) cusp form for $\bG'_A$;

\item the mapping

\[
\pi' \to \pi,
\]

restricted to the collection of (greater than one dimensional) cusp forms on $\bG'_A$ is 1-1 onto the collection of all (equivalence classes of) cusp forms $\bigotimes_v \pi_v$ on $\GL(2, \A)$ such that $\pi_v$ is square-integrable for each $v \in S$.

\end{enumerate}

\end{theorem}

\begin{remark*}

Suppose $\pi'$ has central character $\psi$. Then $\pi'$ is called a cusp form if it occurs in the natural representation $R'_\psi$ defined by right translation in the space $\bL^2(\bG'_F \backslash \bG'_A, \psi)$ consisting of measurable functions $\phi$ on $\bG'_F \backslash \bG'_A$ satisfying

\[
\phi(zg') = \psi(z)\phi(g) \ (\text{for all } z \in Z'_A = Z_A)
\]

and

\[
\int_{Z'_A \bG'_F \backslash \bG'_A} |\phi(g)|^2 \dd g < \infty.
\]

Note that $\pi'$ will occur in the representation $R'$ acting in $\bL^2(\bX')$ only if $\psi$ is trivial on $Z_\infty^+$. More generally (from the definition of $\pi(\pi')$ and the identity immediately following \eqref{eq:7.16}) it follows that $\pi(\pi')$ again induces the character $\psi$ on $Z_A$. Therefore a more accurate statement of Theorem~\ref{thm:10.5} is that

\[
\pi' \to \pi
\]

map the constituents of $R'_\psi$ one-to-one onto the constituents $\otimes \pi_v$ of $R^{\psi}_0$ such that $\pi_v$ is square-integrable for each $v \in S$.

\end{remark*}

One possible proof of Part (i) of Theorem~\ref{thm:10.5} has already been alluded to. By analogy with Tate's theory for $\GL(1)$ of a field one introduces the local zeta functions

\begin{equation}\label{eq:10.4}
\zeta(f, \pi'_v, s) = \int_{G'_v = D_v^{\times}} f(x)\pi'_v(x)|N(x)|_v^{s + 1/2} \dd^{\times}x\tag{10.4}
\end{equation}

corresponding to each Schwartz-Bruhat function $f$ on $D_v$, each irreducible representation $\pi'_v$ of $D_v^{\times}$, and complex number $s$. These integrals converge for $\operatorname{Re}(s)$ sufficiently large and extend to meromorphic functions in all of $\C$. Furthermore there exist Euler factors $L(\pi'_v, s)$ (of degree $\leq 1$) and scalar factors $\varepsilon(\pi'_v, s)$ such that

\begin{equation}\label{eq:10.5}
\frac{\zeta(f, \pi'_v, s)}{L(\pi'_v, s)}\ \varepsilon(\pi'_v, s) = -\frac{\zeta(\tilde{f}, \tilde{\pi}'_v, 1 - s)}{L(\tilde{\pi}'_v, 1 - s)}\tag{10.5}
\end{equation}

(cf.\ \eqref{eq:6.33}). When $\pi'_v$ is of dimension $> 1$ the factor $L(\pi'_v, s) = 1$. More generally,

\begin{equation}\label{eq:10.6}
L(\pi'_v, s) = L(\pi_v(\pi'_v), s)\tag{10.6}
\end{equation}

\[
L(\tilde{\pi}'_v, s) = L(\tilde{\pi}_v(\pi'_v), s)
\]

and

\begin{equation}\label{eq:10.7}
\varepsilon(\pi'_v, s) = \varepsilon(\pi_v(\pi'_v), s)\tag{10.7}
\end{equation}

for all $\pi'_v$. In this sense (among many others) the Weil correspondence $\pi'_v \to \pi_v(\pi'_v)$ is completely natural.

Now still following Tate, one can show that if $\pi'$ is a cusp form then the Euler products

\[
L(s, \pi') = \prod_v L(s, \pi'_v)
\]

and

\[
L(s, \tilde{\pi}') = \prod_v L(s, \tilde{\pi}'_v),
\]

initially defined for $\operatorname{Re}(s)$ sufficiently large, define meromorphic functions which are bounded in vertical strips of finite width and satisfy the functional equation

\[
L(s, \pi') = \varepsilon(s, \pi')L(1 - s, \pi')
\]

if

\[
\varepsilon(s, \pi') = \prod_v \varepsilon(s, \pi'_v).
\]

Furthermore these products define entire functions of $s$ if $\pi'$ is not of the form $g \to \omega(N(g))$ for some grossencharacter $\omega$ of $F$. But for each character $\chi_v$ of $F_v$, $\pi_v(\pi'_v) = \pi_v(\tilde{\pi}'_v)$ and $\pi_v(\pi'_v \otimes \chi_v) = \pi_v(\pi'_v) \otimes \chi_v$ (lifting $\chi_v$ to $D_v$ thru the norm map $N_v$). Consequently by \eqref{eq:10.6}, \eqref{eq:10.7} and Part (ii) of Jacquet-Langlands' Theorem~\ref{thm:6.18}, it follows that $\pi = \otimes \pi_v$ must be a cusp form on $\bG_A$, whence Part (i) of Theorem~\ref{thm:10.5}.

To prove Part (ii) of Theorem~\ref{thm:10.5} it seems (at present) to be necessary to use the Selberg trace formula. Actually the approach we have in mind makes it possible to prove Parts (i) and (ii) together in one fell swoop. This approach is based on the observation that the local correspondence

\[
\pi'_v \to \pi_v
\]

(for each ramified place $v$) is natural from yet another point of view, namely character theory.

More precisely, let $R$ denote a set of representatives $E_v$ of (equivalence classes of) quadratic extensions of $F_v$. For simplicity assume that $F_v$ has characteristic zero so that every such $E_v$ is separable. For each $E_v$ choose an imbedding of $E_v$ in $M(2, F_v)$ and $D_v$ and recall that by the theorem of Skolem-Noether this imbedding is uniquely determined up to inner automorphism.

Now the multiplicative group of each extension $E_v$ imbeds as a Cartan subgroup of $\bG_v (= \GL(2, F_v))$ or $\bG'_v (= D_v^{\times})$. Thus $R$ also indexes the set of (equivalence classes of) Cartan subgroups so obtained. In $\bG'_v$ this set indexes a complete set of representatives of Cartan subgroups but in $\bG_v$ it indexes only the so-called elliptic Cartans. The point is that

\begin{equation}\label{eq:10.8}
\chi_{\pi_v}(b) = -\chi_{\pi'_v}(b)\tag{10.8}
\end{equation}

whenever $b$ is regular in $E_v$ (regular means that the roots of the equation $x^2 - \operatorname{tr}(b)x + N(b) = 0$ are distinct). This miraculous fact comprises the content of Proposition 15.5 of \cite{JacquetLanglands}.

According to \eqref{eq:10.8} then, the ``characters'' of the global representations $\pi$ and $\pi'$ should be the same. Before using this fact to prove Theorem~\ref{thm:10.5} we need to quote without proof a somewhat technical but quite believable result.

Suppose $G$ is any locally compact unimodular group and $Z$ is a closed subgroup of its center. For each character $\psi$ of $Z$ let $\bL^1(\psi)$ denote the algebra (under convolution product) of all measurable $f$ on $G$ which satisfy $f(zg) = \psi^{-1}(z)f(g)$ for all $z \in Z$ and whose absolute values are integrable on $Z \backslash G$. Let $S(G)$ denote some dense subalgebra of $\bL^1(\psi)$ which is closed under the operation $f(g) \to f^*(g) = f(g^{-1})$.

\begin{lemma}\label{lem:10.6}

Suppose $\pi^1$ and $\pi^2$ are two unitary (not necessarily irreducible) representations of $G$ which induce the character $\psi$ on $Z$. Suppose also that the operators

\[
\pi^i(f) = \int_{Z \backslash G} f(g)\pi^i(g) \dd g \quad (i = 1, 2)
\]

are Hilbert-Schmidt for all $f$ in $S(G)$. Then $\pi^1$ and $\pi^2$ are equivalent iff

\begin{equation}\label{eq:10.9}
\operatorname{tr} \pi^1(f * f^*) = \operatorname{tr} \pi^2(f * f^*)\tag{10.9}
\end{equation}

for every $f$ in $S(G)$.

\end{lemma}

Note that since $\pi^i(f * f^*) = \pi^i(f)\pi^i(f)^*$, the operators $\pi^i(f * f^*)$ are of trace class by our assumption on $\pi^i(f)$. The lemma is believable because for an irreducible unitary representation $\pi$ the distribution $\pi(f)$ is known to determine $\pi$.

\subsection*{Proof of Theorem~\ref{thm:10.5}}

\begin{proof}

Suppose first that $\pi'$ is an irreducible unitary representation of $\bG'_A$ of dimension greater than one. So as not to obscure the basic idea of the proof let us also assume that each $\pi'_v$ has dimension greater than one for all $v$. (This is actually an assumption on the dimensionality of $\pi'_v$ for $v$ in $S$ only since for the remaining places it can be shown that $\pi'_v$ is infinite dimensional as soon as $\otimes \pi'_v$ is greater than one dimensional.)

Now let $\pi = \pi(\pi')$ denote the corresponding representation of $\bG_A$ and $\psi$ its central character. What has to be shown is that $\pi'$ is equivalent to some subrepresentation of $R'_\psi$ if and only if $\pi(\pi')$ is equivalent to some subrepresentation of $R^{\psi}_0$. To prove this we shall apply the criterion of Lemma~\ref{lem:10.6}.

To decide which $G$ should be used in Lemma~\ref{lem:10.6} let us consider the decomposition

\[
R^{\psi}_0 = \oplus \pi^j.
\]

Each irreducible constituent $\pi^j = \otimes \pi^j_v$ acts on a certain subspace $V^j$ of $\bL_0^2(\psi)$ so we let $V^j_v$ denote the space of $\pi^j_v$. Next we restrict our attention to the direct sum of those representations $\pi^i$ whose $v$-th components are equivalent to $\pi_v(\pi'_v)$ for each $v \in S$. We denote this subrepresentation of $R^{\psi}_0$ by $R_S$.

For each $v \in S$ we fix a $\bK_v$-finite unit vector $u_v$ in each $V^i_v$ which is the image (under $\pi'_v \to \pi_v(\pi'_v)$) of some fixed $\bK'_v$-finite unit vector $u'_v$ in the space of $\pi'_v$. The direct sum

\[
M = \bigoplus_i \left\{\bigotimes_{v \in S} u_v\right\} \otimes \left\{\bigotimes_{v \notin S} V^i_v\right\},
\]

summed over $i$ indexing $\pi^1$ of the type just described will then be invariant and irreducible under the action of

\[
\bG_S = \{g = (g_v): g_v = 1 \ \text{for all } v \text{ in } S\}.
\]

On the other hand we can consider the decomposition

\[
R'_\psi = \oplus (\pi')^j
\]

and its natural subrepresentation $R'_S$. Replacing $\pi_v(\pi'_v)$ by $\pi'_v$ and $u_v$ by $u'_v$ we can define a subspace $M'$ of $\bL^2(\bX')$ analogous to the subspace $M$ described above. The group

\[
\bG'_S = \{g = (g'_v): g'_v = 1 \ \text{for all } v \in S\}
\]

now acts irreducibly on $M'$ but this group is naturally isomorphic to $\bG_S$ via the local isomorphisms $\theta_v$ described earlier.

Our claim is that to prove Theorem~\ref{thm:10.5} it suffices to prove that these two natural representations of $\bG_S$ (on $M$ and $M'$) are equivalent.

Indeed in $R'_S$ we are dealing with all irreducible unitary representations $\pi'$ of $\bG'_A$ which occur in $R'$ and have $v$-th component equal to $\pi'_v$ for each $v \in S$. The assertion that the corresponding representations $\pi(\pi')$ occur in $R^{\psi}_0$ is equivalent to the assertion that the direct sum of these representations is equivalent to $R_S$. But the $v$-th component of each summand of this direct sum is automatically equivalent to the $v$-th component of each summand of $R_S$ for each $v \in S$. So the last assertion simply says that the remaining components are equal, or, alternately, that the representations of $\bG_S$ in $M$ and $M'$ are equivalent! The argument for the converse direction is entirely similar.

So letting $\tau$ and $\tau'$ denote the representations of $\bG_S$ in $M$ and $M'$ it remains to prove that $\tau$ and $\tau'$ are equivalent. To do this we shall apply Lemma~\ref{lem:10.6} with $G = \bG_S$, $Z = \prod_{v \notin S} Z_v$, $\psi = \prod_{v \notin S} \psi_v$, $\pi^1 = \tau$, $\pi^2 = \tau'$, and $S(G)$ equal to the linear span of the space of functions on $\bG_S$ of the form

\[
f(g) = \prod_{v \notin S} f_v(g_v)
\]

where

\begin{enumerate}[(i)]

\item for all $v$,

\[
f_v(z_v g_v) = \psi_v^{-1}(z_v)f_v(g_v), \quad z_v \in Z_v,
\]

\item for all $v$, $f_v$ is bi-$\bK_v$-finite and of compact support modulo $Z_v$;

\item if $v$ is infinite, $f_v$ is infinitely differentiable;

\item if $v$ is finite, $f_v$ is locally constant; and

\item for almost all finite $v$, $f_v$ is zero outside $Z_v\bK_v$ and on $Z_v\bK_v$ is equal to $\omega_v^{-1}(\det g)$ where $\omega_v$ is unramified and such that $\omega_v^2 = \psi_v$.

\end{enumerate}

According to Lemma~\ref{lem:10.6} what has to be shown is that

\begin{equation}\label{eq:10.10}
\operatorname{tr} \tau(f * f^*) = \operatorname{tr} \tau'(f * f^*)\tag{10.10}
\end{equation}

for all $f$ in $S(G)$. To verify \eqref{eq:10.10} we shall bring into play two natural representations whose traces are easily calculated (using the trace formula) and then relate these traces to those appearing in \eqref{eq:10.10}.

The regular representation relating to the right side of \eqref{eq:10.10} is simply $R'_\psi$. Indeed for each $v$ in $S$ let $f'_v$ denote the function

\[
f'_v(g) = d(\pi'_v)\overline{(\pi'_v(g)u'_v, u'_v)}
\]

where $d(\pi'_v)$ is the formal degree of $\pi'_v$ with respect to (a certain) Haar measure on $\bG'_v$. (The formal degree with respect to $\dd g'_v$ is defined by the orthogonality relations

\begin{equation}\label{eq:10.11}
\int_{Z'_v \backslash \bG'_v} \langle\pi'_v(g'_v)u_1, \tilde{u}_1\rangle\ \overline{\langle\pi'_v(g'_v)u_2, \tilde{u}_2\rangle} \dd g'_v = \frac{1}{d(\pi'_v)}\langle u_1, u_2\rangle\langle\tilde{u}_2, \tilde{u}_1\rangle\tag{10.11}
\end{equation}

for all $u_1, u_2, \tilde{u}_1, \tilde{u}_2$.) If $\Phi'(g)$ is the function on $\bG'_A$ defined by

\[
\Phi'(g) = \left\{\prod_{v \in S} f'_v(g_v)\right\} f * f^*(g'_S)
\]

it follows (from the orthogonality relations \eqref{eq:10.11}!) that $R'_\psi(\Phi')$ coincides with $\tau'(f * f^*)$ on $M'$ and annihilates its orthocomplement. Consequently

\begin{equation}\label{eq:10.12}
\operatorname{trace} R'_\psi(\Phi') = \operatorname{trace} \tau'(f * f^*).\tag{10.12}
\end{equation}

Now let us consider the left side of \eqref{eq:10.10}. Recall that by our assumption on each $\pi'_v$, $v \in S$, the representation $\pi_v(\pi'_v)$ is super cuspidal for such $v$ (assuming $S$ contains no infinite places). Thus the matrix coefficients of $\pi_v(\pi'_v)$ are compactly supported modulo $Z_v$ and in particular satisfy orthogonality relations of type \eqref{eq:10.11} (without the primes everywhere). The measure $\dd g_v$ on $Z_v \backslash \bG_v$ may be normalized so that the formal degree of each $\pi_v(\pi'_v)$ equals $d(\pi_v)$, and consequently we set

\[
f_v(g_v) = d(\pi_v)\overline{(\pi_v(g)u_v, u_v)}
\]

for each $\pi_v = \pi_v(\pi'_v)$, $v \in S$. Defining a function $\Phi(g)$ on $\bG_A$ by

\[
\Phi_f(g) = \left\{\prod_{v \in S} f_v(g_v)\right\} f * f^*(g_S)
\]

if $g = (g_v)g_S$, the result once again is that (with this choice of $\Phi$)

\begin{equation}\label{eq:10.13}
\operatorname{trace} R^{\psi}_0(\Phi) = \operatorname{trace} \tau(f * f^*).\tag{10.13}
\end{equation}

Now according to \eqref{eq:10.12} and \eqref{eq:10.13} we can at last complete the proof of Theorem~\ref{thm:10.5} by establishing that

\[
\operatorname{trace} R^{\psi}_0(\Phi_f) = \operatorname{trace} R'_\psi(\Phi'_f)
\]

for all $f \in S(G_S)$. Not by chance this is precisely the type of identity which the trace formula can establish. Indeed by a slight modification of Remark~\ref{rem:9.23} (replacing $\bL^2(Z_\infty^+ \bG'_F \backslash \bG'_A)$ by $\bL^2(\bG'_F \backslash \bG'_A, \psi)$) we have that

\begin{equation}\label{eq:10.14}
\operatorname{tr} R'_\psi(\Phi'_f) = \sum_{\{\gamma\}} \operatorname{meas}(Z'_A \bG'(\gamma)_F \backslash \bG'(\gamma)_A) \int_{\bG'(\gamma)_A \backslash \bG'_A} \Phi'(x^{-1}\gamma x) \dd x\tag{10.14}
\end{equation}

the summation extending over a set of representatives of conjugacy classes in $Z'_F \backslash \bG'_F$. On the other hand, by Corollary~\ref{cor:9.24}, we have that

\begin{equation}\label{eq:10.15}
\begin{aligned}
\operatorname{tr} R^{\psi}_0(\Phi_f) &= \operatorname{meas}(Z_A \bG_F \backslash \bG_A)\Phi(e) \\
&\quad + \sum_{\substack{\gamma \in \{\bG_F\} \\ (\text{modulo } Z_F)}} \operatorname{meas}(Z_A \bG(\gamma)_F \backslash \bG(\gamma)_A) \int_{\bG(\gamma)_A \backslash \bG_A} \Phi(x^{-1}\gamma x) \dd x
\end{aligned}\tag{10.15}
\end{equation}

the hypothesis of the corollary being satisfied since $S$ contains at least two places (which for convenience we are assuming finite) and for these places the matrix coefficients of $\pi_v(\pi'_v)$ satisfy

\begin{equation}\label{eq:10.16}
\int_{N_v \bK_v} \int f_v(k^{-1}ank) \dd n \dd k = 0.\tag{10.16}
\end{equation}

To complete the proof of Theorem~\ref{thm:10.5} we shall exploit \eqref{eq:10.8} and the discussion immediately surrounding it to show that \eqref{eq:10.14} and \eqref{eq:10.15} are equal.

Assuming for the sake of simplicity that $F$ has characteristic zero, each $\gamma$ in $\bG'_F = D^{\times}$ lies in some separable quadratic extension of $F$ imbeddable in $D$. Thus the conjugacy classes in $Z'_F \backslash \bG'_F$ are indexed by the elements (identified modulo $Z'_F$) of equivalence classes of quadratic extensions $E$ of $F$ which are imbeddable in $D$. Such extensions are precisely those $E$ such that

\begin{equation}\label{eq:10.17}
E \otimes_F F_v \text{ is a field, } v \in S,\tag{10.17}
\end{equation}

and consequently \eqref{eq:10.14} may be rewritten as

\[
\begin{aligned}
\operatorname{trace}(R'_\psi(\Phi'_f)) &= \operatorname{meas}(Z'_A \bG'_F \backslash \bG'_A)\Phi'(e) \\
&\quad + \frac{1}{2}\sum_{Q'} \sum_{\substack{\gamma \in Z'_F \backslash B_F \\ \gamma \notin Z'_F}} \operatorname{meas}(Z'_A B_F \backslash B_A) \int_{B_A \backslash \bG'_A} \Phi'(x^{-1}\gamma x) \dd x,
\end{aligned}
\]

where $Q'$ denotes a set of representatives of equivalence classes of quadratic extensions satisfying \eqref{eq:10.17}, $B_F = B_F(E)$ is the centralizer of $E$ in $\bG'_F$, and $B_A = B_A(E)$ is the centralizer of $E$ in $\bG'_A$. Keeping in mind the definition of $\Phi'_f$ we have finally that

\begin{equation}\label{eq:10.18}
\begin{aligned}
\operatorname{trace} R'_\psi(\Phi'_f) &= \operatorname{meas}(Z'_A \bG'_F \backslash \bG'_A)\left\{\prod_{v \in S} d(\pi'_v)\right\} f * f^*(e) \\
&\quad + \frac{1}{2}\sum_{Q'} \sum_{\substack{\gamma \in Z'_F \backslash B_F \\ \gamma \notin Z'_F}} \operatorname{meas}(Z'_A B_F \backslash B_A)\left\{\prod_{v \in S} \frac{\chi_{\pi'_v}(\gamma^{-1})}{\operatorname{meas}(Z'_v \backslash B_v)}\right\} \int_{B_S \backslash \bG'_S} f * f^*(x^{-1}\gamma x) \dd x
\end{aligned}\tag{10.18}
\end{equation}

Now we deal similarly with \eqref{eq:10.15}. Letting $Q$ denote a set of representatives for the equivalence classes of all quadratic extensions $E$ of $F$, and defining $B_F = B_F(E)$ and $B_A = B_A(E)$ as above, we have that

\[
\begin{aligned}
\operatorname{trace} R^{\psi}_0(\Phi_f) &= \operatorname{meas}(Z_A \bG_F \backslash \bG_A)\left\{\prod_{v \in S} d(\pi_v)\right\} f * f^*(e) \\
&\quad + \frac{1}{2}\sum_Q \sum_{\substack{\gamma \in Z_F \backslash B_F \\ \gamma \notin Z_F}} \operatorname{meas}(Z_A B_F \backslash B_A) \int_{B_A \backslash \bG_A} \Phi(x^{-1}\gamma x) \dd x.
\end{aligned}
\]

We observe that

\[
\int_{B_A \backslash \bG_A} \Phi(x^{-1}\gamma x) \dd x,
\]

is equal to the product of

\begin{equation}\label{eq:10.19}
\prod_{v \in S} \int_{B_v \backslash \bG_v} f_v(x_v^{-1}\gamma x_v) \dd x_v\tag{10.19}
\end{equation}

and

\[
\int_{B_S \backslash \bG_S} f * f^*(x^{-1}\gamma x) \dd x,
\]

where $f_v(g_v)$ is the matrix coefficient $(\pi_v(g_v)u_v, u_v)$ multiplied by the formal degree $d(\pi_v)$. So since $Z = Z'$, $\bG_S = \bG'_S$, and

\[
\operatorname{meas}(Z'_A \bG'_F \backslash \bG'_A) = \operatorname{meas}(Z_A \bG_F \backslash \bG_A)
\]

(provided the Haar measure chosen on $\bG'_A$ is its Tamagawa measure) it remains only to prove that

\begin{equation}\label{eq:10.20}
\int_{B_v \backslash \bG_v} f_v(x_v^{-1}\gamma x_v) \dd x_v = \frac{1}{\operatorname{meas}(Z'_v \backslash B_v)}\chi_{\pi'_v}(\gamma^{-1}),\tag{10.20}
\end{equation}

if $\gamma$ belongs to a subfield of $D$ and equals zero otherwise. But \eqref{eq:10.20} is equivalent (modulo \eqref{eq:10.8}) to the assertion that

\begin{equation}\label{eq:10.21}
\int_{Z_v \backslash \bG_v} (\pi_v(x_v^{-1}\gamma x_v)u_v, u_v) \dd x_v = d(\pi_v)^{-1}\chi_{\pi_v}(\gamma^{-1}),\tag{10.21}
\end{equation}

if $\gamma \in E$ with $E \otimes_F F_v$ a field and

\begin{equation}\label{eq:10.22}
\int_{Z_v \backslash \bG_v} (\pi_v(x_v^{-1}\gamma x_v)u_v, u_v) \dd x_v = 0\tag{10.22}
\end{equation}

otherwise. Both these properties of matrix coefficients of supercuspidal representations are well known (cf.\ Jacquet-Langlands, Section 7, in particular Proposition 7.5; \eqref{eq:10.22} is essentially equivalent to the Condition \eqref{eq:10.16}). Consequently the proof of Theorem~\ref{thm:10.5} is finally complete.

\end{proof}

\begin{remark}\label{rem:10.7}

In the proof of Theorem~\ref{thm:10.5} we (at times implicitly) assumed that:

\begin{enumerate}[(i)]

\item the set of ramified primes in $D$ does not include infinite primes (i.e.\ real places, since no non-trivial division algebra is defined over $\C$); and

\item $\pi'_v$ has dimension greater than one for all ramified $v$ so that $\pi_v(\pi'_v)$ is super cuspidal as opposed to special.

\end{enumerate}

It should be stressed that these assumptions were made purely for convenience of exposition. If $\pi_v(\pi'_v)$ is special for $v$ ramified in $D$ of course one can no longer define $\Phi_f$ using a matrix coefficient of $\pi_v(\pi'_v)$ since such coefficients are not integrable let alone compactly supported modulo $Z_v$. However one can still find a compactly supported function $f_v$ whose Fourier transform (roughly speaking) is zero everywhere except at $\pi_v(\pi'_v)$ (where it is one). So using this $f_v$ to define $\Phi_f$ the trace formula is applicable and one still concludes that

\[
\operatorname{tr} R^{\psi}_0(\Phi_f) = \operatorname{tr} \tau(f * f^*).
\]

In particular assumption (ii) may be dropped.

As for assumption (i), suppose that $v = \infty$ does ramify in $D$ and that $\pi'_v(h)$ is of the form $N(h)^{-n/2}\rho_n(h)$. Then $\pi_v(\pi'_v)$ is equivalent to the discrete series representation $\sigma(\mu_1, \mu_2)$ of ``lowest weight'' $n + 2$ (cf.\ \eqref{eq:7.19}) and in this case, as has already been remarked in the ``Notes and References'' of the last section, there still exists a compactly supported function $f_v$ which does the job (cf.\ \cite{DufloLabesse}, Section II.4). Thus the proof of Theorem~\ref{thm:10.5} is complete.

\end{remark}

\section{Construction of Some Special Automorphic Forms in the Case of Compact Quotient}\label{sec:10c}

To describe at least one automorphic form $\pi'$ on a given quaternion algebra $D$ defined over $\Q$ we shall construct an automorphic form $\pi$ on $\GL(2)$ of the type described in Theorem~\ref{thm:10.5}. This means that the $p$-th component of $\pi$ must be square-integrable for each $p$ dividing the discriminant of $D$. To construct $\pi$ we shall use grossencharacters of a quadratic extension $L$ of $\Q$ since this is the only construction of cusp forms thus far described. The delicate point will be to choose $\lambda$ so that $\pi(\lambda)$ satisfies the requisite properties.

Recall that $\pi(\lambda) = \otimes \pi_p$ is defined as follows if $\lambda = \prod_v \lambda_v$ is a grossencharacter of $L$. In the terminology of Section~\ref{ch:7},

\[
\pi_p = \pi(\lambda_v, \lambda_{v'})
\]

if $p$ splits in $L$ and both $v$ and $v'$ divide $p$. On the other hand, if $p$ does not split in $L$, i.e.\ if $L \otimes_{\Q} \Q_p$ is a field, then $\pi_p$ equals the irreducible component $\pi_p(\lambda_v)$ of the Weil representation attached to the extension $L \otimes_{\Q} \Q_p$.

Recall that $\pi_p(\lambda_v)$ is a principal series representation if $\lambda_v$ is trivial on the norm one group of $L_v$ and a super cuspidal representation otherwise. Therefore the $p$-th component of $\pi(\lambda)$ will be square integrable for each $p \in S$ if and only if

\begin{equation}\label{eq:10.23}
L \otimes_{\Q} \Q_p \text{ is a field}\tag{10.23}
\end{equation}

for all $p \in S$ and $\lambda_v$ does not factor through the norm map $N_v$ for any $v$ dividing $p \in S$.

This discussion implies that we should consider only quadratic extensions $E$ of $\Q$ which are subfields of $D$, and grossencharacters $\lambda = \prod \lambda_v$ of $E$ such that $\lambda_v$ is ``suitable'' in the above sense for all $v$ lying above the primes of $S$. In particular we have:

\begin{theorem}\label{thm:10.8}

Let $L$ denote a quadratic subfield of $D$. Then:

\begin{enumerate}[(i)]

\item To each ``suitable'' character $\lambda$ of the adele group of $L^{\times}$ one can attach an irreducible admissible representation $\pi'(\lambda)$ of $\bG'_A$, the adele group of $D^{\times}$;

\item If $\lambda$ is a grossencharacter of $L$ (i.e.\ an automorphic form for $L$) then $\pi'(\lambda)$ is an automorphic form for $\bG'_A$, i.e.\ $\pi'(\lambda)$ occurs in some $R'_\psi$.

\end{enumerate}

\end{theorem}

\begin{proof}

Theorems \ref{thm:7.11}, \ref{thm:10.5}, and the above remarks.

\end{proof}

\begin{remark}\label{rem:10.9}

Consider the right regular representation of $\SL(2, \R)$ in $\bL^2(\Gamma \backslash \SL(2, \R))$ with $\Gamma$ as in (10.3). Thus $\Gamma$ is the norm one group of a maximal order of a division algebra $D$ and

\[
\Gamma \backslash \SL(2, \R) \cong Z_\infty^+ \bG'_\Q \backslash \bG'_A / \prod_{p < \infty} \bK'_p.
\]

A natural problem is to decompose this representation into irreducibles.

Since $\Gamma \backslash \SL(2, \R)$ is compact we know a priori that the decomposition is discrete. Furthermore each irreducible $\pi$ of $\SL(2, \R)$ occurs at most finitely many times and with multiplicity equal to the dimension of the corresponding space of automorphic forms of type $\pi$. (If $\pi = \pi_k^+$, $f(z)$ is holomorphic in $\operatorname{Im}(z) > 0$ and of weight $k$ with respect to $\Gamma$; if $\pi = \pi_s^+$ or $\pi_s^C$, then $f(z)$ is invariant with respect to $\Gamma$ and satisfies

\[
\Delta f = -y^2\left(\frac{\partial^2}{\partial x^2} + \frac{\partial^2}{\partial y^2}\right)f = \frac{1 - s^2}{4}f.
\])

Now the multiplicity with which $\pi_k^+$ occurs may be computed using the Selberg trace for compact quotient as explained in Section~\ref{ch:9}.A. Although this method fails to give a multiplicity formula for $\pi_s^+$ or $\pi_s^C$ we do know that such representations occur. Indeed since $\bL^2(\Gamma \backslash \SL(2, \R)/\SO(2)) = \bL^2(\Gamma \backslash \{\operatorname{Im}(z) > 0\})$ is obviously infinite-dimensional it follows that infinitely many class 1 representations of $\SL(2, \R)$ occur in $\bL^2(\Gamma \backslash \SL(2, \R))$. (Cf.\ the proof of Theorem~\ref{thm:2.11}.) The problem is which ones?

In the case of non-compact quotient we constructed examples of real analytic cusp forms using Theorem~\ref{thm:7.11}. Since Theorem~\ref{thm:10.8} is an analogue of Theorem~\ref{thm:7.11} (replacing $\GL(2, \A)$ by $\bG'_A$) one should be able to use it to construct similar examples in the case of compact quotient. To explain this let us consider at least one specific example.

To simplify matters fix $D$ to be the unique (indefinite) division quaternion algebra defined over $\Q$ which splits at all primes except $p = 2$ and 3. The quadratic extension

\[
L = \Q(\sqrt{2})
\]

of $\Q$ satisfies \eqref{eq:10.23} at least for $p = 2$ and 3 and consequently imbeds as a maximal subfield of $D$. (It is perhaps the ``simplest'' such subfield since $L \otimes_{\Q} \Q_p$ is ramified only when $p = 2$.)

According to Theorem~\ref{thm:10.8} we should pick a grossencharacter $\prod \lambda_v$ of $L$ such that $\lambda_v$ is ``suitable'' for $p = 2$ and 3 and (say) unramified for all remaining $p$. Note that by our choice of $L$ the infinite prime $p = \infty$ splits so $\lambda_{v_\infty}$ and $\lambda'_{v_\infty}$ determine a continuous series representation $\pi(s_1, s_2)$ of $\GL(2, \R)$. Consequently the construction of such grossencharacters is tantamount to the construction of examples of real analytic cusp forms for compact quotient. However, rather than pursue this point, we prefer to describe some related questions motivated by the remarks below.

\end{remark}

\begin{remark*}

For $\GL(2)$ there is the global ``multiplicity one'' result of Theorem~\ref{thm:5.7} whose proof uses Whittaker models (hence Fourier expansions of cusp forms along $\left\{\begin{pmatrix} 1 & x \\ 0 & 1 \end{pmatrix}\right\}$). Since this proof obviously does not extend to division quaternion algebras it is of interest to ask whether multiplicity one itself still holds.

\end{remark*}

\begin{theorem}\label{thm:10.10}

In the decomposition of $R'_\psi$ an irreducible representation of $\bG'_A$ occurs at most once.

\end{theorem}

\begin{proof}

If $\pi'$ occurs twice then the corresponding $\pi(\pi')$ would have to occur twice in $R^{\psi}_0$ contradicting multiplicity one for $\GL(2)$. Indeed Theorem~\ref{thm:10.5} establishes a one-to-one correspondence between (most of) the constituents of $R'_\psi$ and certain constituents of $R^{\psi}_0$. In particular it establishes an isomorphism between the subspaces of $\bL^2$ in which these representations act.

\end{proof}

\begin{remark}\label{rem:10.11}

There is an important generalization of Theorem~\ref{thm:10.8} which must be true but thus far seems too difficult to prove.

Suppose $D$ is a division algebra of arbitrary rank ($n^2$) over $F$. Then one can define $\bG'_A$ and $\bK'$ and $\bG'_F$ as in the case of rank 4. (Indeed $D_v$ will now be isomorphic to $\GL(m_v, H_v)$ where $H_v$ is a division algebra of rank $n^2/m_v^2$ over $F_v$ and $v$ will be said to ramify in $D$ only if $m_v < n$.) The natural conjecture to make is:

\end{remark}

\begin{conjecture}\label{conj:10.12}

Suppose $L$ is a maximal subfield of $D$ and $\psi$ is a character of $L_A^{\times}$. Then one may attach to $\psi$ an irreducible admissible representation of $\bG'_A$. If $\psi$ is a grossencharacter and ``suitable'' then $\pi'(\psi)$ occurs in the space of automorphic forms for $\bG'_A$.

\end{conjecture}

How might one go about proving this conjecture? Unfortunately the proof of Theorem~\ref{thm:10.8} does not go through since one knows little about $L$-functions for $\GL(n)$. (In particular one does not have a ``Converse'' to Hecke's theory at the present time.) An alternate approach might be to apply the trace formula directly to $G'$ without recourse to $\GL(n)$. Since this approach seems promising let us sketch what it already involves when $n = 2$.

According to Selberg's trace formula

\begin{equation}\label{eq:10.24}
\begin{aligned}
\operatorname{tr}(R'_\psi(f)) &= \int_{Z_A \bG'_F \backslash \bG'_A} \sum_{\gamma \in Z_F \backslash \bG'_F} \operatorname{tr}(f(x^{-1}\gamma x)) \dd x\\
&= \sum_{\gamma \in \{Z_F \backslash \bG'_F\}} \operatorname{meas}(Z_A \bG'(\gamma)_F \backslash \bG'(\gamma)_A) \int_{\bG'(\gamma)_A \backslash \bG'_A} \operatorname{tr}(f(x^{-1}\gamma x)) \dd x\\
&= \sum_{\pi'} \operatorname{tr} \pi'(f)
\end{aligned}\tag{10.24}
\end{equation}

if $R'_\psi = \oplus (m_{\pi'})\pi'$ and $f$ is a nice function on $\bG'_A$ (with values in the space of operators on some $K$-module $\Lambda$; this extension to operator valued functions $f$ is straightforward).

Quite simply the idea is to pick $f = \prod f_v$ so as to isolate $\pi'(\psi)$ in \eqref{eq:10.24} in much the same way one isolates the multiplicity of a discrete series representation at infinity by choosing $f_\infty$ to be a matrix coefficient for this representation.

If $v$ ramifies in $D$ there is no problem since any matrix coefficient of $\pi'_v(\psi)$ will do the job. Now what happens at the remaining places?

At every unramified place which does not split in $L$, $\pi'_v(\psi)$ is either supercuspidal or a principal series representation of the form $\pi(\delta_v, \delta_v\eta_v)$ (with $\eta_v$ equal to the character of order two of $F_v^{\times}$ associated to $L_v$). If $\pi'_v(\psi)$ is supercuspidal again there is no problem but of course this situation arises only finitely often. For the remaining non-splitting places, $\pi_v(\psi)$ is the extension to $\GL(2)$ of the (reducible) principal series representation of $\SL(2)$ induced from the character $\eta_v$ and this representation (like all induced representations) can not easily be isolated by Fourier transform.

What hope there is for completing this approach stems from the fact that each irreducible piece of the restriction of $\pi(\delta_v, \delta_v\eta_v)$ to $\SL(2)$ may be viewed a ``limit'' of discrete series representations. At the infinite prime this concept of ``limit'' is more than mere formal analogy so the guess is that these representations do distinguish themselves among principal series representations (vis-\`a-vis the trace formula).

Supposing for the moment that for all non-splitting $v$ we could find $f_v$ so that

\[
\int_{\bG(\gamma)_v \backslash \bG_v} f_v(x^{-1}\gamma x) \dd x = 0
\]

whenever $\gamma$ lies in some quadratic extension of $F$ which splits at $v$ (the matrix coefficients of supercuspidal representations of $\GL(2, F_v)$ have already been remarked to provide such functions) it would then follow that all conjugacy classes in $Z_F \backslash \bG'_F$ not indexed by elements of the extension $L$ (describing $\pi(\psi)$) would contribute nothing to the trace in \eqref{eq:10.24}. Indeed quadratic extensions of $F$ which don't split at any place where $L$ doesn't split must equal $L$. In short we would have a formula for $\sum m_{\pi'} \operatorname{tr} \pi'(f)$ (summed over all $\pi'$ occurring in $R'_\psi$) which involved only the extension $L$.

Unfortunately the derivation of such a formula must await further developments in harmonic analysis.

\section{Theta Series Attached to Quaternary Quadratic Forms}\label{sec:10d}

In this subsection we describe how a special case of Theorem~\ref{thm:10.5} may be applied to the so-called ``basis problem'' in the classical theory of automorphic forms.

Recall that theta-series attached to positive-definite quadratic forms define automorphic forms for $\Gamma_0(N)$. (Cf.\ Example B(iv) of Section~\ref{ch:1}.)

Roughly speaking, the ``basis problem'' asks whether all elements of $S_k(\Gamma_0(N), \psi)$ may be expressed as linear combinations of such theta-series in particular theta-series attached to the norm form of a (definite) division quaternion algebra. In representation theoretic language this amounts to the explicit construction of the space of $\pi = \otimes \pi_p$ in $\bL_0^2(\bG_\Q \backslash \bG_A, \psi)$ given that $\pi$ corresponds to $\pi'$ through Theorem~\ref{thm:10.5}.

Our exposition here follows \cite{Shimizu}. The work involved is analogous to the construction of Shalika-Tanaka described in Section~\ref{ch:7}C and the analogy is especially strong as regards the use of Weil's representation. In the present context this representation is used to construct (generalized) theta-functions attached to automorphic forms on division quaternion algebras as opposed to quadratic extensions.

The applications we have in mind deal with cusp forms of weight $k$ and character $\psi$ for $\Gamma_0(N)$ with $N$ square free. Thus we assume that $F = \Q$ and that $D$ is definite and of discriminant $N = p_1 \cdots p_r$.

\subsection{Weil Representations and Theta Series}

The global Weil representation $r(s)$ of $\SL(2, \A)$ attached to $D$ is defined locally in terms of the completions $D_p$.

For each $p$ let $r_p$ denote the (local) Weil representation corresponding to $D_p = D \otimes_{\Q} \Q_p$ and its norm form. If $p \in S$, $D_p$ is the unique division quaternion algebra defined over $\Q_p$, and $r_p$ is as described in Section~\ref{ch:7}A (\eqref{eq:7.4}--\eqref{eq:7.6}). On the other hand, if $p \notin S$, $D_p$ is identified with the vector space $V = M(2, \Q_p)$ whose norm form is the usual determinant. In this case the (Weil) representation for $\SL(2, \Q_p)$ is described by \eqref{eq:7.4}--\eqref{eq:7.6} with $\omega(a) \equiv 1$, $y = 1$, and

\[
x^\sigma = \begin{pmatrix} x_1 & x_2 \\ x_3 & x_4 \end{pmatrix}^{\sigma} = \begin{pmatrix} x_4 & -x_2 \\ -x_3 & x_1 \end{pmatrix}.
\]

Let $S_0(D_A)$ be the space spanned by all elements $\prod_p \Phi_p$ where for almost all $p$, $\Phi_p$ is the characteristic function of $\mathfrak{P}_p$. Since $r_p(s_p)$ ($s_p \in \SL(2, \mathcal{O}_p)$) fixes the characteristic function of $\mathfrak{P}_p$ for almost all $p$ the formula

\[
r(s)\Phi = \prod_p r(s_p)\Phi_p
\]

is meaningful and defines the so-called global Weil representation of $\SL(2, \A)$ attached to $D$. This representation extends to $S(D_A)$ in such a way that the mapping $(s, \Phi) \to r(s)\Phi$ from $\SL(2, \A) \times S(D_A)$ into $S(D_A)$ is continuous. It also extends to a unitary representation in $\bL^2(D_A)$ such that \eqref{eq:7.36} and \eqref{eq:7.37} hold with $\gamma = 1$.

The space of generalized theta functions corresponding to the global representation $r(s)$ consists of functions

\begin{equation}\label{eq:10.25}
\Theta(r(s)\Phi) = \sum_{\xi \in D}(r(s)\Phi)(\xi) = \Theta_\Phi(s)\tag{10.25}
\end{equation}

(cf.\ \eqref{eq:7.38} and \eqref{eq:7.39}). The series in \eqref{eq:10.27} converges uniformly and the function of $s$ it defines is continuous and left invariant for $\SL(2, F)$. To construct a subspace of $\bL_0^2(\bG_A, \psi)$ which realizes the representation $\pi(\pi')$ we decompose this space of theta-functions according to the irreducible representations $\pi'$, i.e.\ we project $\Theta_\Phi(s)$ onto a cusp form which transforms under $\bG'_A$ ``according to $\pi'$'' (analogous to the projection \eqref{eq:7.40}).

\subsection{Decomposition of the Weil Representation}

To decompose $r(s)$ we introduce certain spherical functions attached to $\pi'$. By $\bH^1$ we denote the norm one group of $D$ and then we put

\[
\bK_p^1 = \bK'_p \cap \bH_p^1
\]

and

\[
\bK^1 = \bK' \cap \bH_A^1 = \prod \bK'_p.
\]

If $d = \otimes d_p$ is any irreducible representation of $\bK^1$ which occurs in the restriction of $\pi'$ to $\bK^1$ we let $\bL(d)$ denote the space of all $\phi$ in the space of $\pi'$ such that

\begin{equation}\label{eq:10.26}
\int_{\bK^1} \chi_d(k_1^{-1})\pi'(k_1)\phi \dd k_1 = \phi.\tag{10.26}
\end{equation}

(Here $\chi_d(k_1) = (\dim d)$ trace $(d(k_1))$.)

Since $\bL(d)$ is finite-dimensional we can define the spherical function $\omega_d$ of $d$ with respect to $\pi'$ as

\[
\omega_d(g) = \text{trace } (P(d)\pi'(g))
\]

where $P(d)$ denotes projection onto $\bL(d)$. From this definition it follows that

\[
\omega_d(g) = \sum_{i = 1}^{N}(\pi(g)\phi_i, \phi_i)
\]

if $\{\phi_1, \cdots, \phi_N\}$ is an orthonormal basis of $\bL(d)$. For the application we have in mind, the multiplicity of $d$ in $\bL(d)$ will be one, and in this case we have

\begin{equation}\label{eq:10.27}
\omega_d(g) = \frac{\dim(d)}{\phi(g_0)}\int_{\bK^1}\phi(g_0k_1gk_1^{-1}) \dd k_1\tag{10.27}
\end{equation}

for any non-zero $\phi$ in $\bL(d)$ and any $g_0$ in $\bG'_A$ such that $\phi(g_0) \neq 0$. Since $\pi' = \otimes \pi'_v$, we always have

\begin{equation}\label{eq:10.28}
\omega_d(g) = \prod_p \omega_{d_p}(g_p).\tag{10.28}
\end{equation}

Now we can describe the projection (or ``component'') of $\Theta_\Phi(s)$ which ``transforms according to $\pi'$.'' First we restrict attention to $\Phi = \prod \Phi_p$ such that $\Phi(k_1xk_1^{-1}) = \Phi(x)$ for all $k \in \bK^1$,

\begin{equation}\label{eq:10.29}
\int_{\bK^1} \chi_d(k_1)\Phi(k_1x) \dd k_1 = \Phi(x),\tag{10.29}
\end{equation}

and

\begin{equation}\label{eq:10.30}
\Phi_\infty(x) = e^{-2\pi N(x)}P(N(x)) \text{ trace } \rho_n(x)\tag{10.30}
\end{equation}

$P$ being some polynomial. For such $\Phi$, and $s$ in

\[
\GL(2, \A)^+ = \{s \in \GL(2, \A): \det(s) = N(h), \ h \in D_A^{\times}\}
\]

we put

\begin{equation}\label{eq:10.31}
\Theta_\Phi(d, \pi')(s) = |\det(s)|_A \sum_{i = 1}^{N} \int_{Z'_A \bG'_\Q \backslash \bG'_A} \overline{\phi_i(g)} \int_{\bH^1_\Q \backslash \bH^1_A} \phi_i(nhg) \sum_{\xi \in D} r(s_1)\Phi^g(\xi nh) \dd n \dd g\tag{10.31}
\end{equation}

where $N(h) = \det(s)$, $\Phi^g(x) = \Phi(g^{-1}xg)$, and $s = \begin{pmatrix} \det(s) & 0 \\ 0 & 1 \end{pmatrix}s_1$ with $s_1 \in \SL(2, \A)$.

In case $d$ has multiplicity one in $\bL(d)$, \eqref{eq:10.31} should simply read

\begin{equation}\label{eq:10.32}
\Theta_\Phi(d, \pi')(s) = \frac{\dim(d)|\det(s)|}{\phi(g_0)} \int_{\bH^1_\Q \backslash \bH^1_A} \phi(nhg_0)\left(\sum_{\xi \in D} r(s_1)\Phi^{g_0}(\xi nh)\right) \dd n\tag{10.32}
\end{equation}

where $\phi \in \bL(d)$ and $\phi(g_0) \neq 0$. Note that in this case $\Theta_\Phi(d, \pi')$ is defined completely analogously to the theta function $\theta_\Phi(\lambda)$ described by \eqref{eq:7.40} (and attached to the ``spherical function'' $\lambda$). In any case $\Theta_\Phi(d, \pi')$ may be extended to a function on $\GL(2, \A)$ which is left invariant by $\GL(2, \Q)$ and an element of $A_0(\psi)$.

More precisely, let $\varepsilon_-$ denote the element of $\GL(2, \A)$ whose infinite component is $\begin{pmatrix} -1 & 0 \\ 0 & 1 \end{pmatrix}$ and whose remaining components are trivial. Let $V^*(\pi')$ be the subspace of $A_0(\psi)$ spanned by $\Theta_\Phi(d, \pi')$ (for all $\Phi$ of the type described above) and their translates by $\varepsilon_-$. Then $V^*(\pi')$ is invariant for the natural action of the Hecke algebra $\mathcal{H}(\bG_A)$ and equivalent to (i.e.\ realizes) $\pi(\pi')$.

To prove this it is convenient to analyze the first Fourier coefficients

\[
W_\Theta(g) = \int_{\Q \backslash \A} \Theta\left(\begin{pmatrix} 1 & x \\ 0 & 1 \end{pmatrix} g\right)\overline{r(x)} \dd x
\]

of each $\Theta$ in $V^*(\pi')$. Indeed the map

\[
\Theta \to W_\Theta
\]

commutes with the right action of $\bG_A$ so it suffices to show that this image space (call it $W^*(\pi')$) is equivalent to $\pi(\pi')$.

Now if $\Theta = \Theta_\Phi$,

\[
W_\Theta(g) = |\det g|_A \int_{\bH_A^1} \omega_d(nh)r(g_1)\Phi(nh) \dd n
\]

if $g = \begin{pmatrix} \det(g_1) & 0 \\ 0 & 1 \end{pmatrix}g_1$ belongs to $\GL(2, \A)^+$, and equals zero otherwise. But by \eqref{eq:10.28}, $\omega_d(g) = \prod \omega_{d_p}(g_p)$, so since

\[
\Phi = \prod \Phi_p, \quad \text{and} \quad r(x) = \otimes r_p(s)
\]

matters reduce to analyzing the local functions

\[
W_{\Theta_v}(g_v) = |\det g_v| \int_{\bK_v^1} \omega_{d_p}(nh)r_p(g_1^v)\Phi_p(nh) \dd n
\]

on $\GL(2, \Q_p)$. These functions in turn are shown to generate the Whittaker spaces of $\pi_p(\pi'_p)$ and therefore the desired conclusion follows.

\subsection{Application to the Basis Problem}

In this section we fix $k$ and restrict our attention to representations $\pi'$ of $\bG'_A$ satisfying the following conditions:

\begin{enumerate}[(i)]

\item $\pi'$ has central character $\psi$;

\item $\pi'$ is a constituent of $R'_\psi$;

\item $\pi'_\infty(h) = N(h)^{-(k-2)/2}\rho_{k-2}(h)$ (compare the notation of \eqref{eq:7.18});

and

\item the restriction of $\pi'_p$ to $\bK'_p$ contains the identity representation for all finite $p$.

\end{enumerate}

Note that condition (ii) implies that for $p \in S$ we have $\pi'_p = \chi_p \cdot N$ with $\chi_p$ an unramified character of $\Q_p^{\times}$. In particular, $\pi'_p$ is one-dimensional for all $p \in S$. The representations $\pi(\pi') = \otimes \pi_p$ corresponding to these $\pi'$ then satisfy the following properties:

\begin{enumerate}[(i)]

\item $\pi(\pi')$ has central character $\psi$;

\item $\pi(\pi')$ is a constituent of $R^{\psi}_0$ (by Part (i) of Theorem~\ref{thm:10.5});

\item $\pi_\infty = \sigma(\mu_1, \mu_2)$ where

\[
\mu_1(a) = |a|^{r + k - 3/2}
\]

and

\[
\mu_2(a) = |a|^{r - 1/2}[\operatorname{sgn}(a)]^{k - 2}
\]

(cf.\ \eqref{eq:7.19} and the notation of Section~\ref{ch:4}A);

and

\item the restriction of $\pi_p$ to $\bK_p(1) = \left\{\begin{pmatrix} a & b \\ c & d \end{pmatrix} \in \GL(2, \mathcal{O}_p): c \equiv 0\ (p)\right\}$ contains the identity representation.

\end{enumerate}

(This follows from Theorem~\ref{thm:7.6}(ii) and table \eqref{eq:4.20}.)

Now we apply the construction of the last paragraph to these special $\pi'$ and a special choice of $d$, namely

\[
d = d_\infty \otimes (\otimes d_p)
\]

where $d_\infty$ is equivalent to $\rho_{k-2}$ (identifying $\bK_\infty^1$ with $\SU(2)$) and $d_p$ is the identity representation. For this $d$ matters simplify considerably since $d$ has multiplicity one in $\bL(d)$ (and therefore \eqref{eq:10.32} is relevant).

In particular we may introduce the theta-functions

\[
\theta_{i, j}(g) \quad (i, j = 1, \dots, n)
\]

as follows. Let $H'$ denote the space of $\phi$ in the space of any $\pi'$ right invariant for $\prod_{p < \infty} \bK'_p$. By our assumptions above this space is isomorphic to the space of all functions on $\bG_\Q \backslash \bG_A$ taking values in the space $U$ of matrix coefficients of $\pi'_\infty$ satisfying

\[
\phi(hgk) = \left(\bigoplus_{\dim \pi'_\infty} \pi'_\infty\right)(g^{-1})\phi(h)
\]

for all $h \in \bG'_A$, $g \in \bG'_\infty$, and $k \in \prod \bK'_p$. So letting $\{\phi_i\}_{i = 1, \cdots, n}$ denote any basis of the subspace of $H'$ taking values in some fixed irreducible subspace of $U$ we use \eqref{eq:10.32} and put

\[
\theta_{i, j}(g) = \Theta_\Phi(d, \pi')(g)
\]

substituting $\phi_i$ for $\phi$ and $g_j$ for $g_0$ provided $\det(\phi_i(g_j)) \neq 0$; for $\Phi$ we take $\prod \Phi_p$ where $\Phi_\infty(x) = e^{2\pi N(x)}\operatorname{tr}(\rho_{k-2}(x^{\sigma}))$ and $\Phi_p$ is the characteristic function of $\mathfrak{P}_p$ for all finite $p$.

\begin{theorem}\label{thm:10.13}

Let $H(\psi, N, k)$ denote the space of functions $\phi$ in $A_0(\psi)$ satisfying conditions (ii)--(iv) of \eqref{eq:3.6} of Section~\ref{ch:3}. Let $H_0(\psi, N, k)$ denote the intersection of $H(\psi, N, k)$ with those subspaces of $\bL_0^2(\bG_\Q \backslash \bG_A, \psi)$ belonging to any $\pi$ as above. Then $H_0(\psi, N, k)$ is spanned by the collection $\{\theta_{i, j}\}$, $i, j = 1, \cdots, n$.

\end{theorem}

\begin{proof}

The intersection of $H(\psi, N, k)$ with any of these (finitely many) appropriate $\pi$ is one-dimensional and corresponds to a new form in $S_k(N, \psi)$. Suppose $\pi'$ is the automorphic form on $\bG'_A$ such that $\pi(\pi') = \pi$ (Theorem~\ref{thm:10.5}, Part (ii)). Then there is a $(d, \pi')$ ``spherical function'' $\phi_{\pi'}$ such that $\theta_{\phi, g}$ belongs to the space of $\pi(\pi')$ and so since $\{\phi_{ij}\}$ generates all possible $(d, \pi')$ spherical functions (for these $\pi'$) the theorem follows.

\end{proof}

\begin{remark}\label{rem:10.14}

(Concerning the classical content of Theorem 10.12.) The theta-functions $\theta_{ij}$ intersect the theta-series defined classically using quaternary quadratic forms and spherical harmonics (these harmonics being the basic spherical functions for $\SU(2)$).

In \cite{Eichler3} it is proved that $S_k(\Gamma_0(p))$ is spanned by old forms for $\Gamma_0(p)$ and theta-series (with spherical harmonics of weight $k-2$) attached to a definite division quaternion algebra of discriminant $p$. Keeping in mind that elements of $H(\psi, N, k)$ are new forms for $S_k(N, \psi)$ it follows that Theorem 10.12 contains this result.

\end{remark}

\section*{Notes and References}

\addcontentsline{toc}{subsection}{Notes and References}

The arithmetical theory of quaternion algebras is developed in \cite{Eichler} and summarized among other places, in \cite{Eichler2}.

Various Dirichlet series and automorphic forms on quaternion algebras are discussed in \cite{Shimura3}, \cite{Tamagawa}, \cite{Godement4} and \cite{Selberg}. The general theory is described in [Jacquet-Godement] and \cite{JacquetLanglands}.

Application of the trace formula to automorphic forms on quaternion algebras is the subject matter of \cite{Eichler3, Eichler4}, \cite{Shimizu}, \cite{Shimizu2}, and \cite{JacquetLanglands}. Our exposition is adapted from \cite{JacquetLanglands}. Further arithmetical results relating to the trace formula are described in \cite{Mautner2} and \cite{Langlands6}.
